\documentclass{abntex2}
\usepackage[utf8]{inputenc}
\usepackage[T1]{fontenc}
\usepackage{geometry}
\usepackage{graphicx}
\usepackage{float}
\usepackage{hyperref}
\usepackage[english,indonesian]{babel}

\title{Optimasi Pemrosesan Template LaTeX Menggunakan Pendekatan Berbasis AI}
\author{Academic Support Agent}
\date{\today}

\begin{document}
\maketitle

\begin{abstract}

\end{abstract}

\chapter{PENDAHULUAN}


\chapter{TINJAUAN PUSTAKA}


% Original content (commented out):
% \documentclass[
	% -- opções da classe memoir --
	12pt,				% tamanho da fonte
	openright,			% capítulos começam em pág ímpar (insere página vazia caso preciso)
	oneside,			% para impressão em recto e verso. Oposto a oneside
	a4paper,			% tamanho do papel. 
	% -- opções da classe abntex2 --
	chapter=TITLE,		% títulos de capítulos convertidos em letras maiúsculas
	%section=TITLE,		% títulos de seções convertidos em letras maiúsculas
	%subsection=TITLE,	% títulos de subseções convertidos em letras maiúsculas
	%subsubsection=TITLE,% títulos de subsubseções convertidos em letras maiúsculas
	% -- opções do pacote babel --
	english,			% idioma adicional para hifenização
	brazil,				% o último idioma é o principal do documento
	]{abntex2}


% ------------------------------------
% PACOTES
% ------------------------------------

% Pacotes fundamentais 
\usepackage{lmodern}			% Usa a fonte Latin Modern
\usepackage[T1]{fontenc}		% Selecao de codigos de fonte.
\usepackage[utf8]{inputenc}		% Codificacao do documento (conversão automática dos acentos)
\usepackage{indentfirst}		% Indenta o primeiro parágrafo de cada seção.
\usepackage{color}				% Controle das cores
\usepackage{graphicx}			% Inclusão de gráficos
\usepackage{microtype} 			% para melhorias de justificação
\usepackage{IFMG_Arcos} 		% formato específico relatório IFMG Campus Arcos 
% Pacotes adicionais, usados no anexo do modelo de folha de identificação
\usepackage{multicol}
\usepackage{multirow}
	
% Pacotes adicionais, usados apenas no âmbito do Modelo Canônico do abnteX2
\usepackage{lipsum}				% para geração de dummy text

% Pacotes de citações
\usepackage[brazilian,hyperpageref]{backref}	 % Paginas com as citações na bibl
\usepackage[alf]{abntex2cite}	% Citações padrão ABNT

% ------------------------------------
% CONFIGURAÇÕES DE PACOTES
% ------------------------------------

% Configurações do pacote backref
% Usado sem a opção hyperpageref de backref
\renewcommand{\backrefpagesname}{Citado na(s) página(s):~}
% Texto padrão antes do número das páginas
\renewcommand{\backref}{}
% Define os textos da citação
\renewcommand*{\backrefalt}[4]{
	\ifcase #1 %
		Nenhuma citação no texto.%
	\or
		Citado na página #2.%
	\else
		Citado #1 vezes nas páginas #2.%
	\fi}%

% Informações de dados para CAPA
\titulo{Relatório TAI \\ }
% ---------------------------------------------------
%Inserir nomes dos autores aqui
\autor{	Autor 1 \\
		Autor 2 \\
		Autor 3 \\
		Autor 4 \\
		Autor 5}
% ---------------------------------------------------
\local{Arcos - MG}
\data{Dezembro - 2019}
\tipotrabalho{Relatório TAI}
% ---------------------------------------------------
%Inserir nomes dos autores aqui
% O preambulo deve conter o tipo do trabalho, o objetivo, 
% o nome da instituição e a área de concentração 
\preambulo{Relatório apresentado ao Instituto Federal de Educação, Ciência e Tecnologia de Minas Gerais (IFMG), referente ao Trabalho Acadêmico Integrador, como requisito parcial para aprovação em todas as disciplinas ministradas no 6$^\circ$ período do Curso de Engenharia Mecânica.}



% Configurações de aparência do PDF final

% alterando o aspecto da cor azul
\definecolor{black}{RGB}{0,0,0}

% informações do PDF
\makeatletter
\hypersetup{
     	%pagebackref=true,
		pdftitle={\@title}, 
		pdfauthor={\@author},
    	pdfsubject={\imprimirpreambulo},
	    pdfcreator={LaTeX with abnTeX2},
		pdfkeywords={abnt}{latex}{abntex}{abntex2}{relatório técnico}, 
		colorlinks=true,       % false: boxed links; true: colored links
    	linkcolor=black,       % color of internal links
    	citecolor=black,       % color of links to bibliography
    	filecolor=black,       % color of file links
		urlcolor=black,
		bookmarksdepth=4
}
\makeatother

% Espaçamentos entre linhas e parágrafos 

% O tamanho do parágrafo é dado por:
\setlength{\parindent}{1.3cm}

% Controle do espaçamento entre um parágrafo e outro:
\setlength{\parskip}{0.2cm}  % tente também \onelineskip

% compila o indice
\makeindex

% Início do documento
\begin{document}

% Seleciona o idioma do documento (conforme pacotes do babel)
%\selectlanguage{english}
\selectlanguage{brazil}

% Retira espaço extra obsoleto entre as frases.
\frenchspacing 

% ----------------------------------------------------------
% ELEMENTOS PRÉ-TEXTUAIS
% ----------------------------------------------------------

% Capa
\imprimircapa

% Folha de rosto
% (o * indica que haverá a ficha bibliográfica)
\imprimirfolhaderosto*

% Dedicatória

\begin{dedicatoria}
   \vspace*{\fill}
   \centering
   \noindent
   \textit{...} \vspace*{\fill}
\end{dedicatoria}

% Agradecimentos

\begin{agradecimentos}
	\vspace{\fill}
	


\end{agradecimentos}

% Epígrafe
\begin{epigrafe}
    \vspace*{\fill}
	\begin{flushright}
		\textit{"As pessoas pensam que sonhos não são reais \\ apenas porque não são feitos de matéria, de partículas.\\ Sonhos são reais. Mas eles são feitos de \\ pontos de vista, de imagens, de memórias e trocadilhos, \\ e de esperanças perdidas."  Neil Gaiman}
	\end{flushright}
\end{epigrafe}

% resumo em português
\setlength{\absparsep}{18pt} % ajusta o espaçamento dos parágrafos do resumo
\begin{resumo}



\textbf{Palavras-chave}:
\end{resumo}

% resumo em inglês
\begin{resumo}[Abstract]
 \begin{otherlanguage*}{english}
   
   
   \vspace{\onelineskip}
 
 
 
   \noindent 
   \textbf{Key words}:  
 \end{otherlanguage*}
\end{resumo}

%inserir lista de ilustrações
\pdfbookmark[0]{\listfigurename}{lof}
\listoffigures*
\cleardoublepage

% inserir lista de tabelas
% ---
\pdfbookmark[0]{\listtablename}{lot}
\listoftables*
\cleardoublepage



% inserir o sumario
\pdfbookmark[0]{\contentsname}{toc}
\tableofcontents*
\cleardoublepage

\input{textos/resumo.tex}

% ----------------------------------------------------------
% ELEMENTOS TEXTUAIS
% ----------------------------------------------------------
\textual

\input{textos/introducao.tex}

\input{textos/referencial_teorico.tex}

\input{textos/metodologia.tex}

\input{textos/desenvolvimento.tex}

\input{textos/conclusoes.tex}

% ----------------------------------------------------------
% ELEMENTOS PÓS-TEXTUAIS
% ----------------------------------------------------------
\postextual

% ----------------------------------------------------------
% Referências bibliográficas
% ----------------------------------------------------------
\bibliography{referencias}

% ----------------------------------------------------------
% Apêndice
% ----------------------------------------------------------
\begin{apendicesenv}
\input{textos/apendice.tex}

\end{apendicesenv}

% ----------------------------------------------------------
% Anexo
% ----------------------------------------------------------

\begin{anexosenv}
\input{textos/anexo.tex}

\end{anexosenv}


\end{document}


\end{document}